\documentclass[11pt,aspectratio=169]{beamer}
\usetheme{metropolis}
\usepackage{booktabs}
\usepackage{graphicx}
\usepackage{amsmath}

\title{The Asymmetric Timeliness of Earnings Revisited}
\subtitle{A replication of Basu (1997) on U.S.\ firms, 2015--2024}
\author{Justin Cheng}
\institute{University of Oklahoma}
\date{ECON 5253, Spring 2026}

\begin{document}

\maketitle

% ---- 1. Introduction ----
\begin{frame}{Topic and motivation}
  \textbf{Question.} Do U.S.\ public firms still book bad news into earnings
  faster than good news, the way Basu (1997) reported for 1963--1990?
  \medskip
  \begin{itemize}
    \item Conservatism principle: gains require more verifiable evidence
          than losses, so earnings should react more strongly to negative
          returns.
    \item The Basu reverse regression became the workhorse test of
          conditional conservatism.
    \item Critics (Dietrich--Muller--Riedl, 2007) and modern fair-value
          rules leave it open whether the asymmetry survives in the
          post-2015 sample.
  \end{itemize}
\end{frame}

% ---- 2. Data ----
\begin{frame}{Data}
  \textbf{Source.} WRDS Compustat \texttt{funda} + CRSP \texttt{msf},
  linked through the CCM link table.

  \medskip
  \textbf{Sample filters.}
  \begin{itemize}
    \item U.S.\ common stocks (CRSP \texttt{shrcd} $\in \{10,11\}$),
          NYSE, AMEX, NASDAQ.
    \item December fiscal year-ends. Drop financials (SIC 6000--6999).
    \item Beginning-of-period price above one dollar; no missing EPS or
          price.
    \item Winsorize $EP$ and $R$ at the 1\% / 99\% pooled tails.
  \end{itemize}

  \medskip
  \textbf{Variables.}
  $EP_{it} = \mathrm{EPS}_{it} / P_{i,t-1}$ (split-adjusted),
  $R_{it}$ is the compounded twelve-month return,
  $D_{it} = \mathbf{1}\{R_{it} < 0\}$.
\end{frame}

% ---- 3. Methods ----
\begin{frame}{Methods}
  \textbf{Basu (1997) reverse regression}
  \begin{equation*}
    EP_{it} = \beta_0 + \beta_1 D_{it} + \beta_2 R_{it}
            + \boxed{\beta_3} \, D_{it} R_{it} + \varepsilon_{it}
  \end{equation*}
  \begin{itemize}
    \item $\beta_2$: sensitivity of earnings to good news.
    \item $\beta_2 + \beta_3$: sensitivity of earnings to bad news.
    \item $\beta_3 > 0$ means conditional conservatism is present.
  \end{itemize}

  \medskip
  Estimated by OLS, with firm-clustered standard errors (Petersen 2009).
  Sub-samples: pre-COVID (2015--2019) vs.\ COVID era (2020--2024).
  Year-by-year coefficients drive the trend figure.
\end{frame}

% ---- 4. Findings ----
\begin{frame}{Findings}
  \begin{itemize}
    \item Pooled sample: $\beta_3 = 0.402$ ($t=26.4$). Conservatism is
          alive in 17{,}188 firm-years.
    \item Implied bad-news slope $\beta_2 + \beta_3 \approx 0.32$;
          the good-news slope $\beta_2 \approx -0.08$ is near zero.
    \item Pre- and COVID-era averages: $0.450 \to 0.364$. Looks like
          mild compression on its own.
    \item Year-by-year tells the real story.
          \begin{itemize}
            \item 2020 jumps to $0.60$ as COVID impairments hit.
            \item 2021 collapses to $0.13$ in the recovery year, since
                  the prior write-downs cannot reverse.
            \item By 2023 and 2024 the coefficient is back near $0.5$.
          \end{itemize}
  \end{itemize}
\end{frame}

% ---- 5. Takeaway ----
\begin{frame}{What did we learn?}
  \begin{itemize}
    \item Basu (1997) replicates on a fresh decade of U.S.\ data. Bad
          news loads on earnings far faster than good news.
    \item Conditional conservatism is state-dependent. The COVID-19
          shock pushed $\beta_3$ up in 2020, the recovery pulled it
          down in 2021, and two-period averages hide that swing.
    \item Practical takeaway: earnings remain a lopsided signal.
          Analysts, lenders, and standard-setters should not assume the
          asymmetry has gone away under modern fair-value rules.
    \item Limitations: the Dietrich--Muller--Riedl (2007) mechanical-bias
          critique still applies; the standard errors here are
          firm-clustered only, with two-way clustering left for future
          work.
  \end{itemize}
\end{frame}

\end{document}
